% ============================================================
% Come si usano, argomento 10: variabili aleatorie continue
% ============================================================
\argomento{Variabili aleatorie continue: densità e CDF}{Formule a p.~\pageref*{f:puntozero}}[Densità e CDF]

\begin{scelta}{Come scegliere la formula}
\renewcommand{\arraystretch}{1.25}
\begin{tabular}{@{}P{0.55\linewidth}P{0.39\linewidth}@{}}
  \toprule
  \textbf{Se il testo\ldots} & \textbf{si usa\ldots} \\
  \midrule
  chiede la probabilità di un intervallo & l'integrale della densità \vv{intervallo}, o la differenza della CDF
    \vv{cdf} \\
  chiede la probabilità di un singolo valore & è zero \vv{puntozero} \\
  chiede se una funzione è una densità, o la costante che la rende tale & \vv{proprietapdf}: positiva, area 1 \\
  dà la densità e chiede la CDF, o viceversa & integrale fino a $x$, oppure derivata \vv{cdf} \\
  chiede «$X$ supera $x$» & la CCDF, $1 - F_X(x)$ \vv{cdf} \\
  chiede il valore atteso & \vv{mediacontinua} \\
  \bottomrule
\end{tabular}

Nel continuo tutto si fa come nel discreto, con l'integrale al posto della somma e $f_X(x) \dd x$ al posto di
$p_X(x)$. Se la densità è definita a pezzi, si spezzano gli integrali sugli stessi pezzi.
\end{scelta}

\scheda{puntozero}{
  \usoquando{Si chiede la probabilità di un valore esatto di una variabile continua, o bisogna decidere se gli estremi
    di un intervallo contano.}
  \usoattenzione{Probabilità zero non vuol dire impossibile: ogni misura dà \emph{qualche} valore, e quel valore aveva
    probabilità zero.}
  \usoesempio{Per una variabile continua $P(X \leq 2) = P(X < 2)$, e $P(1 \leq X \leq 3) = P(1 < X < 3)$.}
}

\scheda{pdf}{
  \usoquando{Per descrivere una variabile continua: la densità dice dove la variabile cade più spesso (dove è alta)
    e dove raramente (dove è bassa).}
  \usocome{Per un intervallo piccolo la probabilità è circa la densità per l'ampiezza dell'intervallo.}
  \usoattenzione{$f_X(x)$ non è una probabilità e può superare 1: la probabilità è l'area, non l'altezza.}
  \usoesempio{Uniforme su $\left[0, \frac{1}{2}\right]$: $f_X = 2$, e $P(0{,}10 \leq X \leq 0{,}11) \simeq 2 \cdot
    0{,}01 = 0{,}02$.}
}

\scheda{intervallo}{
  \usoquando{Si chiede la probabilità che $X$ cada fra due valori, o sopra o sotto una soglia.}
  \usocome{Si integra la densità fra gli estremi, trovando una primitiva (\vv{fondcalcolo}). Se la densità è definita a
    pezzi, si spezza l'integrale; fuori dal supporto la densità è zero e non contribuisce.}
  \usoesempio{$f_X(x) = 2x$ su $[0, 1]$: $P\left(X \leq \frac{1}{2}\right) = \int_0^{1/2} 2x \dd x =
    \left[x^2\right]_0^{1/2} = \frac{1}{4}$.}
}

\scheda{proprietapdf}{
  \usoquando{Si deve controllare che una funzione sia una densità, o trovare la costante $c$ che la rende tale.}
  \usocome{Si controlla che sia $\geq 0$ ovunque; poi si impone che l'area totale valga 1 e si ricava $c$.}
  \usoesempio{$f_X(x) = c\, x^2$ su $[0, 1]$: $\int_0^1 c\, x^2 \dd x = \frac{c}{3} = 1$, quindi $c = 3$.}
}

\scheda{cdf}{
  \usoquando{Si chiede «$X$ non supera $x$» (CDF) o «$X$ supera $x$» (CCDF); oppure molte probabilità di intervalli,
    che con la CDF diventano semplici differenze; oppure si deve passare dalla densità alla CDF o viceversa.}
  \usocome{Dalla densità alla CDF: si integra da $-\infty$ fino a $x$, un pezzo alla volta, aggiungendo l'area
    accumulata nei pezzi precedenti. Dalla CDF alla densità: si deriva.}
  \usoesempio{Densità triangolare $1 - \abs{x}$ su $[-1, 1]$: $F_X(x) = \frac{(1 + x)^2}{2}$ per $-1 \leq x \leq 0$ e
    $F_X(x) = 1 - \frac{(1 - x)^2}{2}$ per $0 \leq x \leq 1$. Quindi $P\left(-\frac{1}{2} \leq X \leq \frac{1}{2}\right)
    = \frac{7}{8} - \frac{1}{8} = \frac{3}{4}$.}
}

\scheda{proprietacdf}{
  \usoquando{Per controllare che una funzione sia una CDF, o che una CDF calcolata sia giusta: deve crescere da 0 a 1
    senza mai scendere.}
  \usovarianti{Per una variabile discreta la CDF è una scala: piatta fra un valore e il successivo, con un salto alto
    $p_X(x)$ in ogni valore $x$.}
  \usoesempio{Un dado onesto: $F_X(3{,}5) = P(X \leq 3) = \frac{1}{2}$, e la CDF sale di $\frac{1}{6}$ in ciascuno dei
    valori $1, \ldots, 6$.}
}

\scheda{mediacontinua}{
  \usoquando{Si vuole il valore atteso di una variabile continua: il baricentro della densità.}
  \usocome{Si integra $x\, f_X(x)$ sul supporto. Se la densità è simmetrica attorno a un punto $c$ e la media esiste,
    la media è $c$ senza fare conti.}
  \usoattenzione{L'integrale deve convergere: la variabile di Cauchy (argomento 11) non ha media.}
  \usoesempio{$f_X(x) = 2x$ su $[0, 1]$: $E[X] = \int_0^1 2x^2 \dd x = \frac{2}{3}$, spostata a destra del punto di
    mezzo. La densità triangolare, simmetrica attorno a 0, ha media 0.}
}
